\section{Introduction}
\label{sec:introduction}

Electrifying bus fleets requires coordinated decisions on charging infrastructure, battery capacity, and daily charging operations.
Wireless charging technologies expand the set of feasible charging opportunities by allowing buses to charge at selected stops or road links during regular service.
This creates a planning trade-off: more in-route wireless charging can reduce battery requirements and residual plug-in fast-charging needs, but installing such infrastructure is costly and must be robust to future service plans.

The planning difficulty is amplified by service-plan uncertainty.
Transit agencies may adjust routes, timetables, and duty assignments over a medium-term planning horizon, while embedded wireless facilities are shared physical investments that cannot be changed after the scenario is realized.
Thus, a first-stage wireless charging layout must support multiple possible representative operating days.

Wireless infrastructure also changes downstream operational coupling.
After a service-plan scenario is realized, the operator chooses route-level battery capacities and charging schedules.
SWC and DWC opportunities affect how much energy can be obtained during service, which in turn changes battery sizing and the residual demand for existing plug-in fast chargers.
Because plug-in fast chargers are limited shared operational resources, residual fast-charging demand couples otherwise route-level recourse decisions.

These features lead to a two-stage stochastic mixed-integer planning problem.
The first stage selects shared wireless charging facilities, and the second stage contains scenario-dependent battery and charging-operation decisions.
The recourse problem has discrete battery choices, operational charging decisions, and shared plug-in fast-charger capacity constraints, making the deterministic equivalent and a direct integer L-shaped decomposition computationally challenging.

This paper develops an exact reformulation-enabled integer L-shaped method for this problem class.
The three algorithmic contributions are nested rather than independent.
First, we introduce an exact structure-exposing lifted reformulation that places scenario-dependent battery configuration variables in the master while preserving their scenario-adaptive role.
Second, we use the lifted structure to build a nested scenario--route relaxation--verification decomposition: Phase 1 relaxes shared plug-in fast-charger capacity coupling and operational integrality to obtain route-level LP lower bounds, and Phase 2 restores the exact scenario-level mixed-integer operational recourse.
Third, we derive joint-resource direction-selective monotonicity cuts that exploit the partial order induced by wireless facility additions and battery-capacity upgrades.
The dependency chain is lifting, route-decomposable relaxation, exposed joint-resource order, and strengthened cuts.

The computational study uses Tower-derived electric-bus planning instances as a realistic test bed for the algorithm.
The experiments first give a compact planning interpretation, then evaluate the algorithm through network-scale ablation, scenario-count scalability, and capacity-coupling stress tests.
The evidence is intended to show which components are necessary for scalability and how the method behaves when the Phase 1 relaxation becomes more or less aligned with the exact Phase 2 feasible region.
